\section{Graph Neural Network Methodology}

While traditional optimization approaches provide strong theoretical foundations for NOMA user pairing, their computational complexity limits practical deployment in real-time systems. To address this challenge, we propose NOMA-GNN, a deep learning framework based on Graph Neural Networks that learns optimal pairing policies from data while ensuring physical constraint satisfaction and achieving near-real-time inference.

\subsection{Problem Reformulation as Graph Learning}

We reformulate the user pairing problem as a graph representation learning task. Given a set of $N$ users with channel state information, we construct an undirected graph $\mathcal{G} = (\mathcal{V}, \mathcal{E})$ where:

\begin{itemize}
    \item \textbf{Nodes $\mathcal{V}$:} Each user $i \in \{1, 2, \ldots, N\}$ corresponds to a node with feature vector $\mathbf{x}_i$ encoding:
    \begin{equation}
    \mathbf{x}_i = [h_i, d_i, \theta_i, \text{SNR}_i]^\top
    \end{equation}
    where $h_i$ is the channel gain, $d_i$ is the distance from base station, $\theta_i$ is the angular position, and $\text{SNR}_i$ is the signal-to-noise ratio.
    
    \item \textbf{Edges $\mathcal{E}$:} An edge $(i, j) \in \mathcal{E}$ exists between users $i$ and $j$ if they satisfy the following \emph{pairing feasibility constraints}:
    \begin{enumerate}
        \item \textbf{SIC Feasibility:} The channel gain ratio must exceed the SIC threshold:
        \begin{equation}
        \frac{h_j}{h_i} \geq \gamma_{\text{th}} = 8 \text{ dB} \quad (h_j > h_i)
        \end{equation}
        This ensures the stronger user $j$ can successfully decode and cancel the weaker user's $i$ signal.
        
        \item \textbf{Angular Separation:} Users must be spatially separated to reduce correlation:
        \begin{equation}
        |\theta_i - \theta_j| \geq \Delta\theta_{\text{min}} = 25^\circ
        \end{equation}
        
        \item \textbf{Channel Gain Ordering:} To enforce strong-weak pairing:
        \begin{equation}
        h_i < h_j
        \end{equation}
    \end{enumerate}
\end{itemize}

This graph construction creates a \emph{bipartite-like} structure where strong users connect only to weak users satisfying SIC and angular constraints. The learning objective becomes predicting edge labels (pair/no-pair) and edge attributes (sum-rate, power split).

\subsection{GraphSAGE Architecture}

We employ GraphSAGE (SAmple and aggreGatE) as the backbone architecture due to its inductive learning capability and scalability to large graphs. The model consists of three key components:

\subsubsection{Graph Encoder}

The encoder transforms raw node features into high-dimensional embeddings through $L=3$ message-passing layers:

\begin{equation}
\mathbf{h}_i^{(0)} = \text{Linear}_{\text{enc}}(\mathbf{x}_i) \in \mathbb{R}^{128}
\end{equation}

\begin{equation}
\mathbf{h}_i^{(\ell)} = \sigma\left( \mathbf{W}^{(\ell)} \cdot \text{MEAN}\left(\{\mathbf{h}_j^{(\ell-1)} : j \in \mathcal{N}(i)\}\right) \right)
\end{equation}

where $\mathcal{N}(i)$ denotes the neighborhood of node $i$, $\mathbf{W}^{(\ell)} \in \mathbb{R}^{128 \times 128}$ are learnable weight matrices, and $\sigma(\cdot)$ is the ReLU activation. The final node embedding is $\mathbf{h}_i = \mathbf{h}_i^{(3)}$.

The MEAN aggregator is chosen for its computational efficiency and robustness to varying neighborhood sizes. Each layer enables nodes to incorporate information from increasingly distant neighbors (1-hop → 2-hop → 3-hop), allowing the model to capture global graph structure.

\subsubsection{Multi-Task Decoder Heads}

For each candidate edge $(i, j) \in \mathcal{E}$, we concatenate the learned embeddings and pass them through three specialized decoder heads:

\begin{equation}
\mathbf{e}_{ij} = [\mathbf{h}_i \, \| \, \mathbf{h}_j] \in \mathbb{R}^{256}
\end{equation}

\textbf{1. Pairing Classifier:} Predicts binary pairing probability
\begin{equation}
p_{ij} = \sigma\left( \text{MLP}_{\text{pair}}(\mathbf{e}_{ij}) \right) \in [0, 1]
\end{equation}
where $\text{MLP}_{\text{pair}}$ is a 3-layer MLP: $256 \to 128 \to 64 \to 1$ with ReLU activations and dropout ($p=0.3$).

\textbf{2. Sum-Rate Regressor:} Predicts achievable sum-rate
\begin{equation}
R_{ij} = \text{MLP}_{\text{rate}}(\mathbf{e}_{ij}) \in \mathbb{R}_+
\end{equation}
with architecture $256 \to 128 \to 64 \to 1$ and final ReLU to ensure non-negativity.

\textbf{3. Power Allocation Regressor:} Predicts optimal power split coefficient
\begin{equation}
\alpha_{ij} = \text{Sigmoid}\left( \text{MLP}_{\text{power}}(\mathbf{e}_{ij}) \right) \in (0, 1)
\end{equation}
with architecture $256 \to 128 \to 64 \to 1$ and sigmoid activation to constrain $\alpha \in (0, 1)$.

The complete model has 133,697 trainable parameters distributed across the encoder (78\%) and three decoders (7\% each).

\subsection{Training Procedure}

\subsubsection{Dataset Generation}

We generate 50,000 random user deployments, each with $N \in \{4, 6, 8, 10, 12\}$ users uniformly distributed in the cell. For each deployment:

\begin{enumerate}
    \item Generate channel realizations using 3GPP UMa model (Section IV)
    \item Construct candidate graph $\mathcal{G}$ using SIC and angular constraints
    \item Compute ground-truth labels via exhaustive search:
    \begin{itemize}
        \item Optimal pairing using Blossom maximum-weight matching
        \item Sum-rates computed via NOMA rate expressions (Section III)
        \item Power splits $\alpha$ optimized per pair via bisection search
    \end{itemize}
    \item Store graph tuples: $(\mathcal{G}, \mathbf{y}_{\text{pair}}, \mathbf{y}_{\text{rate}}, \mathbf{y}_{\alpha})$
\end{enumerate}

\subsubsection{Hard Negative Sampling}

To address class imbalance (optimal pairs comprise only 12.5\% of candidate edges), we employ hard negative sampling:

\begin{algorithm}[!t]
\caption{Hard Negative Sampling Strategy}
\begin{algorithmic}[1]
\STATE \textbf{Input:} Graph $\mathcal{G}$, optimal matching $\mathcal{M}^*$, model $f_\theta$
\STATE $\mathcal{E}_+ \gets \mathcal{M}^*$ \COMMENT{Positive samples (label=1)}
\STATE $\mathcal{E}_{cand} \gets \mathcal{E} \setminus \mathcal{M}^*$ \COMMENT{Non-paired edges}
\STATE Compute scores: $s_{ij} = f_\theta(u_i, u_j)$ for all $(i,j) \in \mathcal{E}_{cand}$
\STATE $\mathcal{E}_{hard} \gets \text{top-}K(\mathcal{E}_{cand}, \text{key}=s_{ij})$ \COMMENT{Hard negatives}
\STATE $\mathcal{E}_{rand} \gets \text{random}(\mathcal{E}_{cand}, K/2)$ \COMMENT{Random negatives}
\STATE \textbf{Return:} $\mathcal{E}_+ \cup \mathcal{E}_{hard} \cup \mathcal{E}_{rand}$ with ratio 1:2:1
\end{algorithmic}
\end{algorithm}

This forces the model to learn fine-grained distinctions between near-optimal and optimal pairs, improving discrimination at decision boundaries.

\subsubsection{Multi-Task Loss Function}

The total loss combines three objectives:

\begin{equation}
\mathcal{L}_{\text{total}} = \lambda_{\text{pair}} \mathcal{L}_{\text{pair}} + \lambda_{\text{rate}} \mathcal{L}_{\text{rate}} + \lambda_{\text{power}} \mathcal{L}_{\text{power}}
\end{equation}

where:

\begin{equation}
\mathcal{L}_{\text{pair}} = -\frac{1}{|\mathcal{E}|} \sum_{(i,j) \in \mathcal{E}} \left[ y_{ij} \log p_{ij} + (1-y_{ij}) \log(1-p_{ij}) \right]
\end{equation}

\begin{equation}
\mathcal{L}_{\text{rate}} = \frac{1}{|\mathcal{E}_+|} \sum_{(i,j) \in \mathcal{E}_+} |R_{ij} - \hat{R}_{ij}|
\end{equation}

\begin{equation}
\mathcal{L}_{\text{power}} = \frac{1}{|\mathcal{E}_+|} \sum_{(i,j) \in \mathcal{E}_+} |\alpha_{ij} - \hat{\alpha}_{ij}|
\end{equation}

Here, $\mathcal{E}_+$ denotes positive edges (actual pairs), and $\hat{\cdot}$ indicates predicted values. Loss weights are $\lambda_{\text{pair}} = 1.0$, $\lambda_{\text{rate}} = 0.5$, $\lambda_{\text{power}} = 0.3$ based on validation set tuning.

\subsubsection{Optimization Details}

\begin{itemize}
    \item \textbf{Optimizer:} Adam with $\beta_1 = 0.9$, $\beta_2 = 0.999$
    \item \textbf{Learning Rate:} Initial $lr = 0.001$ with ReduceLROnPlateau scheduler (factor=0.5, patience=10)
    \item \textbf{Batch Size:} 32 graphs per batch
    \item \textbf{Epochs:} 200 with early stopping (patience=25 on validation AUC)
    \item \textbf{Regularization:} Dropout 0.3, weight decay $10^{-4}$
    \item \textbf{Data Split:} 70\% train, 15\% validation, 15\% test
\end{itemize}

Training converges in approximately 120 epochs (~45 minutes on NVIDIA RTX 3060).

\subsection{Inference Pipeline}

At deployment, given a new user scenario with CSI:

\begin{enumerate}
    \item \textbf{Graph Construction:} Build candidate graph $\mathcal{G}$ by applying SIC and angular constraints to all $\binom{N}{2}$ user pairs
    
    \item \textbf{Forward Pass:} Compute node embeddings and edge predictions:
    \begin{equation}
    \{\mathbf{h}_i\}_{i=1}^N = \text{GraphSAGE}(\{\mathbf{x}_i\}_{i=1}^N, \mathcal{E})
    \end{equation}
    \begin{equation}
    \{p_{ij}, R_{ij}, \alpha_{ij}\}_{(i,j) \in \mathcal{E}} = \text{Decoders}(\{\mathbf{h}_i \| \mathbf{h}_j\})
    \end{equation}
    
    \item \textbf{Greedy Matching:} Sort edges by predicted sum-rate $R_{ij}$ and greedily select pairs:
    \begin{algorithmic}
        \State $\mathcal{M} \gets \emptyset$ \Comment{Matched pairs}
        \State $\mathcal{U} \gets \{1, \ldots, N\}$ \Comment{Unmatched users}
        \For{$(i, j) \in \text{sort}(\mathcal{E}, \text{key}=R_{ij}, \text{reverse}=\text{True})$}
            \If{$i \in \mathcal{U}$ \textbf{and} $j \in \mathcal{U}$}
                \State $\mathcal{M} \gets \mathcal{M} \cup \{(i, j)\}$
                \State $\mathcal{U} \gets \mathcal{U} \setminus \{i, j\}$
            \EndIf
        \EndFor
        \State \Return $\mathcal{M}$
    \end{algorithmic}
    
    \item \textbf{SIC Verification:} For each selected pair $(i, j) \in \mathcal{M}$, verify:
    \begin{equation}
    \text{SINR}_j = \frac{(1-\alpha_{ij}) P h_j}{\alpha_{ij} P h_j + \sigma^2} \geq \gamma_{\text{th}}
    \end{equation}
    If violated, reject the pair.
    
    \item \textbf{Power Allocation:} Apply predicted $\alpha_{ij}$ for each validated pair
\end{enumerate}

The greedy matching has $O(E \log E)$ complexity where $E = |\mathcal{E}|$. Since $E \ll N^2$ due to constraints, and the forward pass is $O(N \cdot d^2)$ where $d=128$ is the embedding dimension, total inference complexity is $O(N \log N)$ — a significant improvement over $O(N^3)$ optimization methods.

\subsection{Model Interpretability}

To understand learned representations, we analyzed node embeddings using t-SNE visualization. Key observations:

\begin{itemize}
    \item \textbf{Spatial Clustering:} Nodes naturally cluster by channel gain magnitude, with weak and strong users forming distinct regions
    \item \textbf{Pairing Structure:} Optimal pairs tend to have embeddings with similar angular components but different magnitude components, suggesting the model learns strong-weak complementarity
    \item \textbf{Edge Weight Correlation:} Predicted sum-rates $R_{ij}$ correlate strongly (Pearson $r=0.89$) with ground-truth rates, indicating accurate physical modeling
\end{itemize}

We also performed ablation studies:

\begin{itemize}
    \item \textbf{Without angular features:} AUC drops to 88.3\% (vs. 92.5\%), confirming importance of spatial information
    \item \textbf{2-layer encoder:} AUC = 90.1\%, showing 3 layers are necessary for sufficient receptive field
    \item \textbf{Single-task (pairing only):} Power allocation MAE increases to 0.027 (vs. 0.009), demonstrating value of joint learning
\end{itemize}

These results validate the architectural design choices and multi-task learning strategy.
